\chapter{Orientasi: Enterprise Architecture, Platform, dan Penulisan Akademik}
\label{ch:orientasi}

\section*{Tujuan Pembelajaran}
\begin{itemize}
  \item Menjelaskan posisi \emph{Enterprise Architecture} (EA) sebagai jembatan
        antara strategi organisasi dan kapabilitas teknologi, dengan bahasa
        manajerial yang dapat dipahami audiens non-IT.
  \item Membedakan empat domain arsitektur enterprise---Business, Data,
        Application, Technology (BDAT)---serta peran API dan platform sebagai
        pengungkit nilai pada model bisnis digital modern.
  \item Mengenali anatomi makalah ilmiah yang relevan untuk bidang Enterprise
        Architecture dan platform/API, serta kriteria memilih target publikasi
        seperti jurnal SINTA, jurnal internasional, atau konferensi akademik.
  \item Menyusun pernyataan topik awal, memilih organisasi studi kasus dalam
        konteks Indonesia, dan merumuskan dua sampai tiga kandidat target
        publikasi sebagai dasar kerja satu semester.
\end{itemize}

\section{Pendahuluan}

Bab ini menjadi pintu masuk untuk seluruh perjalanan mata kuliah. Sebelum
memasuki konsep-konsep seperti TOGAF-lite, arsitektur BDAT, strategi API, dan
tata kelola platform, pembaca diperkenalkan terlebih dahulu pada
\emph{gambaran besar} yang mengikat semua konsep tersebut: bagaimana sebuah
organisasi menerjemahkan strategi menjadi keputusan arsitektur, dan bagaimana
mahasiswa magister non-IT dapat mengolah proses analisis itu menjadi makalah
ilmiah yang layak dipublikasikan.

Tiga gagasan utama bab ini saling berkaitan. Pertama, Enterprise Architecture
diperkenalkan bukan sebagai produk diagram yang rumit, melainkan sebagai cara
berpikir untuk membaca organisasi sebagai sebuah sistem. Kedua, ekonomi
platform dan strategi API diperkenalkan sebagai konteks bisnis modern yang
membuat keputusan arsitektur semakin strategis dan tidak lagi semata-mata
menjadi urusan tim TI. Ketiga, makalah publikasi diperkenalkan sebagai sarana
untuk mengubah proses analisis selama satu semester menjadi kontribusi
akademik yang dapat ditelaah oleh komunitas.

Bagi mahasiswa berlatar belakang non-IT, bab ini menegaskan bahwa keputusan
arsitektur enterprise pada dasarnya adalah keputusan manajerial yang dapat
dipahami melalui contoh bisnis sehari-hari. Bagi mahasiswa yang sudah pernah
bersentuhan dengan teknologi, bab ini menggeser perhatian dari kode dan
implementasi menuju pertanyaan yang lebih strategis: mengapa, untuk siapa,
dengan kompromi apa, dan berdasarkan bukti apa.

\section{Skenario Pembuka: Bank Daerah dan Rencana Transformasi Digital}

Bayangkan sebuah bank daerah di Indonesia, sebut saja \emph{Bank Wanua}, yang
selama tiga puluh tahun melayani nasabah lokal dengan jaringan cabang
tradisional. Direksi baru ingin membawa Bank Wanua masuk lebih serius ke
layanan digital: aplikasi mobile, integrasi pembayaran QRIS, kemitraan dengan
fintech,
serta layanan API publik untuk pelaku UMKM yang menjadi tulang punggung
ekonomi daerah. Dewan direksi sepakat tentang \emph{tujuan}, tetapi mereka
tidak yakin tentang \emph{urutan} dan \emph{prioritas}.

Direktur Operasional bertanya: ``Sistem mana yang harus diperbarui terlebih
dahulu, dan layanan mana yang sebaiknya dibuka sebagai API?'' Direktur
Kepatuhan bertanya: ``Bagaimana Bank Wanua memenuhi UU Perlindungan Data
Pribadi (UU PDP), Standar Nasional Open API Pembayaran (SNAP BI) Bank
Indonesia, dan ketentuan Otoritas Jasa Keuangan (OJK)?'' Direktur Pemasaran
bertanya: ``Apakah Bank Wanua sebaiknya menjadi platform mandiri atau
bergabung sebagai mitra ke ekosistem yang lebih besar?''

Tidak satu pun pertanyaan itu dapat dijawab hanya dengan keputusan teknis.
Setiap jawaban membutuhkan pemahaman tentang strategi bisnis, kapabilitas yang
sudah dimiliki, kondisi data dan aplikasi saat ini, aturan main regulasi
Indonesia, serta dinamika ekosistem mitra. Pertanyaan seperti inilah yang
menjadi wilayah kerja \emph{Enterprise Architecture}.

\subsection*{Pertanyaan Pemandu Bab Ini}

Sepanjang bab, lima pertanyaan berikut menjadi kompas pembaca:

\begin{enumerate}
  \item Apa sesungguhnya yang dimaksud dengan Enterprise Architecture, dan
        mengapa pembaca non-IT perlu memahaminya?
  \item Bagaimana platform dan API mengubah cara organisasi mengambil keputusan
        arsitektur dalam dekade terakhir?
  \item Apa yang membedakan sebuah \emph{makalah ilmiah} dari laporan proyek
        atau dokumen konsultasi?
  \item Bagaimana cara memilih target publikasi yang realistis dan relevan,
        khususnya dalam konteks Indonesia?
  \item Bagaimana pembaca dapat memulai perjalanan satu semester ini dengan
        memilih organisasi studi kasus dan topik makalah yang tepat?
\end{enumerate}

\section{Mengapa Topik Ini Penting bagi Pembaca Non-IT}

Mata kuliah ini sengaja dirancang untuk mahasiswa magister berlatar belakang
non-IT karena keputusan arsitektur enterprise pada akhirnya adalah keputusan
manajerial. Di dalamnya ada pertimbangan bisnis, anggaran, regulasi, sumber
daya manusia, budaya organisasi, dan risiko. Pengembang dan arsitek teknis
dapat membantu pelaksanaan, tetapi arah, prioritas, tata kelola, dan kompromi
yang diambil tetap perlu dipahami oleh pengambil keputusan.

Salah satu kesalahan yang sering terjadi adalah memperlakukan Enterprise
Architecture sebagai ``urusan tim TI''. Cara pandang ini membuat strategi
digital mudah gagal. Penyebabnya bukan selalu teknologi yang buruk, melainkan
keputusan teknis yang tidak tersambung dengan keputusan bisnis. Akibatnya,
organisasi dapat membayar mahal untuk sistem canggih yang jarang digunakan,
atau membangun kapabilitas TI yang tidak relevan dengan pelanggan dan model
bisnisnya.

Dengan memahami EA sebagai bahasa lintas fungsi, pembaca non-IT dapat menjadi
jembatan antara dewan direksi, tim bisnis, tim TI, dan mitra eksternal.
Kemampuan ini penting bagi peran strategis seperti Chief Information Officer,
Chief Digital Officer, Strategy Manager, Product Manager, Business Analyst,
atau pemimpin transformasi digital di sektor publik dan swasta.

\section{Kerangka Konseptual: Empat Pilar Bab Ini}

Empat konsep inti diperkenalkan pada bab pengantar ini:

\begin{enumerate}
  \item \textbf{Enterprise Architecture (EA)} sebagai jembatan strategi dan
        teknologi.
  \item \textbf{BDAT---Business, Data, Application, Technology} sebagai empat
        domain arsitektur yang akan dipetakan secara bertahap.
  \item \textbf{API dan Platform} sebagai pengungkit nilai pada model bisnis
        digital modern.
  \item \textbf{Anatomi Makalah Ilmiah} sebagai bentuk komunikasi akademik yang
        akan dibangun selama 14 sesi.
\end{enumerate}

\subsection{Enterprise Architecture: Jembatan Strategi dan Teknologi}

Enterprise Architecture adalah disiplin yang membantu organisasi menyelaraskan
strategi bisnis dengan kapabilitas teknologi melalui pemetaan, prinsip, dan
keputusan yang terdokumentasi. Hasil EA bukan sekadar kumpulan diagram,
melainkan keputusan terstruktur tentang apa yang perlu dipertahankan,
diperbaiki, diganti, atau dibangun dari awal.

Gambar~\ref{fig:ea-jembatan} menggambarkan posisi EA secara konseptual. EA
menjembatani \emph{lapisan strategi} (visi, model bisnis, dan tujuan
organisasi) dengan \emph{lapisan teknologi} (sistem, data, aplikasi, dan
infrastruktur). Tanpa jembatan ini, kedua sisi sering berbicara dengan bahasa
yang berbeda.

\begin{figure}[htbp]
\centering
\scalebox{0.95}{
\begin{tikzpicture}[
  font=\small\bfseries,
  node distance=12mm,
  layer/.style={
    draw, rounded corners=2mm, align=center,
    minimum width=88mm, minimum height=14mm
  },
  strategi/.style={layer, fill=praditagreen!20},
  ea/.style={layer, fill=orange!25, minimum height=18mm},
  tek/.style={layer, fill=blue!12},
  arrow/.style={Latex-Latex, very thick}
]
\node[strategi] (s) {STRATEGI ORGANISASI\\
\footnotesize\mdseries Visi, model bisnis, tujuan, pelanggan, mitra, regulasi};
\node[ea, below=of s] (e) {ENTERPRISE ARCHITECTURE\\
\footnotesize\mdseries Prinsip, peta kapabilitas, kondisi saat ini$\to$target, gap, peta jalan};
\node[tek, below=of e] (t) {KAPABILITAS TEKNOLOGI\\
\footnotesize\mdseries Aplikasi, data, infrastruktur, integrasi, API, platform};

\draw[arrow] (s.south) -- (e.north);
\draw[arrow] (e.south) -- (t.north);

\node[right=0.3cm of e, align=left, font=\footnotesize\mdseries\itshape, text width=30mm]
 {menerjemahkan\\strategi$\leftrightarrow$\\ kapabilitas};
\end{tikzpicture}
}
\caption{Enterprise Architecture sebagai jembatan dua arah antara strategi
organisasi dan kapabilitas teknologi.}
\label{fig:ea-jembatan}
\end{figure}

Analogi yang dekat adalah cetak biru bangunan. Cetak biru bukanlah bangunan
itu sendiri; ia adalah gambar kerja yang membantu pemilik, kontraktor, pekerja
lapangan, dan pemeriksa berbicara tentang rancangan yang sama. Di dalamnya ada
pilihan tentang jumlah lantai, posisi tangga darurat, material pondasi, sistem
pendingin, urutan pembangunan, dan biaya. Cetak biru juga memuat asumsi dan
prinsip: apakah bangunan tahan gempa, hemat energi, dan ramah disabilitas.
Dengan cara serupa, Enterprise Architecture memuat keputusan dan prinsip
tentang ``bangunan organisasi'' yang sedang dirancang, dirawat, atau
ditransformasikan.

\subsection{BDAT: Empat Domain Arsitektur Enterprise}

Empat domain klasik dalam Enterprise Architecture, yang sering disingkat BDAT,
membantu kita memecah kerumitan organisasi menjadi beberapa lapisan yang lebih
mudah dibaca. Setiap lapisan dapat dianalisis secara terpisah, tetapi
keputusannya tetap saling memengaruhi. Gambar~\ref{fig:bdat-stack}
menunjukkan empat lapisan tersebut.

\begin{figure}[htbp]
\centering
\scalebox{0.95}{
\begin{tikzpicture}[
  font=\small\bfseries,
  node distance=2mm,
  layer/.style={
    draw, rounded corners=1.5mm, align=left,
    minimum width=110mm, minimum height=12mm,
    inner xsep=4mm
  },
  business/.style={layer, fill=praditagreen!25},
  data/.style={layer, fill=orange!20},
  app/.style={layer, fill=blue!15},
  tech/.style={layer, fill=gray!18}
]
\node[business] (b) {\textbf{B} -- BUSINESS\\
\footnotesize\mdseries Kapabilitas, value stream, proses, pemangku kepentingan, model operasional};
\node[data, below=of b] (d) {\textbf{D} -- DATA\\
\footnotesize\mdseries Entitas data inti, kepemilikan, kualitas, privasi, tata kelola};
\node[app, below=of d] (a) {\textbf{A} -- APPLICATION\\
\footnotesize\mdseries Portofolio aplikasi, integrasi, pola layanan, bangun sendiri vs.\ beli};
\node[tech, below=of a] (t) {\textbf{T} -- TECHNOLOGY\\
\footnotesize\mdseries Cloud, infrastruktur, keamanan, prinsip teknologi};
\end{tikzpicture}
}
\caption{Empat domain arsitektur enterprise (BDAT) tersusun dari lapisan
strategi-bisnis di puncak hingga teknologi pendukung di dasar.}
\label{fig:bdat-stack}
\end{figure}

Domain \textbf{Business} menjawab pertanyaan: apa yang dilakukan organisasi
untuk pelanggannya, dan kapabilitas apa yang dibutuhkan untuk melakukannya.
Domain \textbf{Data} menjawab: informasi apa yang dibutuhkan, siapa
pemiliknya, dan bagaimana informasi itu dilindungi. Domain
\textbf{Application} menjawab: aplikasi dan layanan apa yang mendukung
kapabilitas tersebut, serta bagaimana aplikasi saling terhubung. Domain
\textbf{Technology} menjawab: infrastruktur, keamanan, dan platform teknis apa
yang menopang semuanya.

Empat domain ini akan dipetakan secara bertahap pada Bab~5 sampai Bab~8. Pada
bab pengantar ini, pembaca cukup memahami satu hal: hampir setiap keputusan
organisasi akan menyentuh keempat domain tersebut, baik disadari maupun tidak.
Tugas arsitektur enterprise adalah membuat hubungan itu terlihat, sehingga
keputusan dapat dijelaskan dan dipertanggungjawabkan.

Sepanjang buku ini, domain BDAT dinyatakan secara lebih formal menggunakan
\emph{ArchiMate}, bahasa pemodelan standar untuk arsitektur enterprise.
ArchiMate menata arsitektur ke dalam beberapa lapisan: \emph{motivasi} dan
\emph{strategi} di puncak (mengapa dan dengan cara apa), lalu lapisan
\emph{bisnis}, \emph{aplikasi}, serta \emph{teknologi dan fisik} (apa dan
bagaimana), dan akhirnya \emph{implementasi dan migrasi}. Karena itu, sebelum
membedah domain BDAT, Bab~3 dan Bab~4 lebih dahulu menegaskan lapisan motivasi
dan lapisan strategi, baru Bab~5 masuk ke lapisan bisnis. Pemetaan lapisan ke
bab ditampilkan pada Tabel~\ref{tab:peta-lapisan} di bagian Pendahuluan.

\subsection{API dan Platform: Pengungkit Nilai Digital}

Dalam dua dekade terakhir, banyak model bisnis digital tumbuh bukan melalui
alur linear sederhana (produsen $\to$ konsumen), melainkan melalui platform
yang menghubungkan banyak pihak. Gojek menghubungkan pengemudi, penumpang,
mitra warung, mitra layanan makanan, dan pengiklan. Tokopedia menghubungkan
pedagang dan pembeli. BCA Open API dan BRI BRIAPI menghubungkan bank dengan
pengembang fintech yang membangun layanan pembayaran. Platform semacam ini
menghasilkan \emph{network effect}: semakin banyak peserta di satu sisi,
semakin besar pula nilai bagi sisi lainnya.

Application Programming Interface, atau API, dapat dipahami sebagai kontrak
digital yang membuat satu sistem dapat berinteraksi dengan sistem lain secara
terstandar. Bagi pembaca non-IT, API dapat dianalogikan sebagai stop kontak:
bentuk dan tegangannya disepakati, sehingga berbagai perangkat dapat terhubung
tanpa harus dinegosiasikan ulang satu per satu. API yang dirancang dengan baik
membuat integrasi dengan mitra menjadi lebih cepat, lebih konsisten, dan lebih
mudah dikelola.

Dalam konteks Bank Wanua, API dan platform menentukan apakah bank tersebut
akan tetap menjadi penyedia layanan yang tertutup, atau berkembang menjadi
simpul ekosistem yang terhubung dengan UMKM, fintech, dan pemerintah daerah.
Keputusan ini bukan sekadar urusan teknis. Ia menyangkut model bisnis, tata
kelola, manajemen risiko, dan posisi strategis bank dalam ekosistem.

\subsection{Anatomi Makalah Ilmiah}

Bab ini juga memperkenalkan luaran utama mata kuliah: makalah publikasi
individu. Makalah ilmiah berbeda dari laporan proyek karena ia memiliki
struktur yang memungkinkan pembaca akademik memahami masalah, menilai metode,
memeriksa argumen, dan menggunakan kontribusinya. Gambar~\ref{fig:anatomi-makalah}
menampilkan struktur minimum yang akan dibangun sesi demi sesi.

\begin{figure}[htbp]
\centering
\scalebox{0.85}{
\begin{tikzpicture}[
  font=\footnotesize\bfseries,
  node distance=2mm,
  sec/.style={
    draw, rounded corners=1.5mm, align=left,
    minimum width=98mm, minimum height=8mm,
    fill=praditagreen!10,
    inner xsep=3mm
  },
  highlight/.style={sec, fill=orange!20}
]
\node[sec] (a) {Judul, Abstrak (150--250 kata), Kata Kunci (4--6)};
\node[sec, below=of a] (b) {Pendahuluan: latar belakang, masalah, RQ, kontribusi};
\node[sec, below=of b] (c) {Tinjauan Pustaka \& Kerangka Konseptual};
\node[sec, below=of c] (d) {Metode/Pendekatan Analisis (TOGAF-lite atau lainnya)};
\node[sec, below=of d] (e) {Konteks Organisasi/Studi Kasus};
\node[highlight, below=of e] (f) {Analisis Kondisi Saat Ini};
\node[highlight, below=of f] (g) {Rancangan Arsitektur Target (BDAT)};
\node[highlight, below=of g] (h) {Strategi API \& Strategi Platform};
\node[highlight, below=of h] (i) {Tata Kelola Platform, Risiko, Metrik};
\node[highlight, below=of i] (j) {Analisis Kesenjangan, Peta Jalan, Implikasi Manajerial};
\node[sec, below=of j] (k) {Diskusi, Keterbatasan, Arah Lanjutan};
\node[sec, below=of k] (l) {Kesimpulan \& Rekomendasi};
\node[sec, below=of l] (m) {Daftar Pustaka, Lampiran Artefak};
\end{tikzpicture}
}
\caption{Anatomi makalah publikasi yang akan dibangun sesi demi sesi. Bagian
berlatar oranye merupakan bagian teknis-substantif yang menjadi inti
kontribusi dan paling banyak dibangun pada Sesi~6 sampai Sesi~12.}
\label{fig:anatomi-makalah}
\end{figure}

Setiap bagian makalah memiliki tujuan komunikasi tersendiri.
\emph{Pendahuluan} meyakinkan pembaca bahwa masalah penting dan layak
dipelajari. \emph{Tinjauan pustaka} menempatkan makalah dalam percakapan
akademik yang sudah ada. \emph{Metode} menjelaskan bagaimana analisis
dilakukan agar dapat dievaluasi atau direplikasi. \emph{Analisis dan
rancangan} menyajikan kontribusi inti. \emph{Diskusi} mengubah temuan menjadi
makna dan implikasi. \emph{Kesimpulan} menutup argumen dengan jelas dan
menunjukkan arah lanjutan.

\section{Konteks Indonesia: Mengapa Lokal Itu Penting}

Banyak buku Enterprise Architecture ditulis dengan asumsi pasar, regulasi, dan
budaya organisasi Amerika Utara atau Eropa. Kondisi tersebut tidak selalu
sejalan dengan Indonesia. Karena itu, mata kuliah ini menempatkan konteks
Indonesia sebagai bagian inti analisis, bukan sekadar catatan tambahan.

Beberapa regulasi dan standar yang akan sering muncul dalam buku ini antara
lain:

\begin{itemize}
  \item \textbf{UU No.~27/2022 tentang Pelindungan Data Pribadi (UU PDP)}
        membentuk batas etis dan hukum tentang bagaimana data pelanggan boleh
        dikumpulkan, diproses, dan dipertukarkan.
  \item \textbf{UU No.~11/2008 jo.~UU No.~19/2016 tentang Informasi dan
        Transaksi Elektronik (UU ITE)} mengatur tanggung jawab penyelenggara
        sistem elektronik.
  \item \textbf{PP No.~71/2019 tentang Penyelenggaraan Sistem dan Transaksi
        Elektronik} mengatur penyelenggaraan sistem elektronik, termasuk
        aspek kedaulatan data dan klasifikasi
        sistem strategis.
  \item \textbf{Permenkominfo PSE} mewajibkan pendaftaran penyelenggara sistem
        elektronik baik publik maupun privat.
  \item \textbf{Standar Nasional Open API Pembayaran (SNAP BI)} dari Bank
        Indonesia membakukan integrasi API pada layanan pembayaran.
  \item \textbf{Ketentuan Otoritas Jasa Keuangan (OJK)} dan \textbf{Badan Siber
        dan Sandi Negara (BSSN)} memberi rambu untuk sektor keuangan dan
        keamanan siber nasional.
\end{itemize}

Selain regulasi, mahasiswa juga akan belajar dari ekosistem digital Indonesia:
Gojek, Tokopedia, Dana, Halodoc, Ruangguru, BCA Open API, BRI BRIAPI, Telkom
Indonesia, Kementerian PANRB melalui program Satu Data Indonesia, serta
koperasi/BPR dan BUMN di sektor energi, transportasi, dan telekomunikasi.
Studi kasus lokal diperkenalkan secara bertahap agar pembahasan tidak berhenti
pada abstraksi.

Tantangan implementasi yang khas Indonesia juga akan dibahas secara terbuka:
keterbatasan anggaran TI di banyak organisasi publik, kapasitas SDM yang
beragam, kematangan vendor lokal yang tidak merata, infrastruktur digital
yang berbeda antarwilayah, regulasi sektoral yang berlapis, serta budaya
organisasi yang sering kali masih hierarkis.

\section{Studi Kasus Mini: Memilih Posisi Bank Wanua}

Mari kita lanjutkan skenario Bank Wanua. Setelah dewan direksi sepakat untuk
melakukan transformasi digital, sebuah lokakarya satu hari diadakan untuk
menjawab tiga pertanyaan strategis berikut:

\begin{enumerate}
  \item Apakah Bank Wanua akan menjadi platform mandiri yang mengundang mitra
        UMKM dan fintech untuk masuk, atau akan bergabung sebagai mitra ke
        platform yang lebih besar (misalnya melalui SNAP BI dan kemitraan
        dengan dompet digital nasional)?
  \item Layanan apa yang sebaiknya dibuka sebagai \emph{public API} terlebih
        dahulu agar memberi nilai segera bagi UMKM, tanpa membahayakan keamanan
        dan kepatuhan?
  \item Kapabilitas internal apa yang harus diperkuat sebelum API publik
        tersebut layak diluncurkan?
\end{enumerate}

Lokakarya seperti ini, bila hanya diisi diskusi bebas, sering berakhir dengan
opini yang saling bertabrakan tanpa keputusan yang jelas. Enterprise
Architecture menyediakan kerangka analisis untuk memetakan kapabilitas, value
stream, kondisi saat ini, target, kesenjangan, dan peta jalan. Dengan begitu,
diskusi dapat bergerak dari adu pendapat menjadi argumen berbasis bukti.

Buku ini mengikuti keputusan tersebut secara konkret. Dari portofolio
kapabilitas yang disepakati Bank Wanua, \textbf{API Transfer dan Perpindahan
Dana}---``rel'' perpindahan uang yang membungkus BI-FAST---dipilih sebagai
kapabilitas sasaran yang dibedah lapis demi lapis sepanjang Bab~3 sampai Bab~5
(lihat Bagian~\ref{sec:studi-kasus-berjalan} di Pendahuluan). Pembaca menerapkan
pola analisis yang sama pada organisasi pilihannya sendiri.

Pada akhir Bab~13, Bank Wanua (atau organisasi pilihan pembaca) akan memiliki
peta kapabilitas bisnis dengan prioritas yang jelas, peta data beserta
kepemilikan dan klasifikasi sensitivitasnya, portofolio aplikasi dengan
rekomendasi tindak lanjut, rancangan arsitektur teknologi target, portofolio
API yang diprioritaskan berdasarkan dampak dan kelayakan, model platform
dengan analisis \emph{network effect}, piagam tata kelola platform yang
mempertimbangkan regulasi Indonesia, analisis kesenjangan, serta peta jalan
12--24 bulan. Semua artefak itu kemudian dirangkai menjadi satu makalah
publikasi.

\section{Memilih Target Publikasi}

Salah satu keputusan paling penting di Sesi~1 adalah memilih dua atau tiga
target publikasi. Target ini akan menentukan templat, panjang naskah, gaya
sitasi, dan cara argumen disusun. Target yang dipilih sebaiknya realistis
terhadap data dan akses yang dimiliki, sekaligus relevan dengan tema makalah.

\subsection{Jenis Target Publikasi dan Karakteristiknya}

Gambar~\ref{fig:venue-map} merangkum jenis target publikasi yang umum diakses oleh
mahasiswa pascasarjana di Indonesia.

\begin{figure}[htbp]
\centering
\resizebox{0.72\textwidth}{!}{
\begin{tikzpicture}[
  font=\small\bfseries,
  level distance=14mm,
  sibling distance=22mm,
  node distance=10mm,
  block/.style={
    draw, rounded corners=1.5mm, align=center,
    minimum width=30mm, minimum height=10mm
  },
  root/.style={block, fill=praditagreen!25, minimum width=70mm},
  cat/.style={block, fill=orange!20},
  ven/.style={block, fill=blue!10, minimum width=22mm, font=\footnotesize\bfseries}
]
\node[root] (root) {Target Publikasi};
\node[cat, below left=of root, xshift=-10mm] (jur) {Jurnal Nasional\\(SINTA 2--4)};
\node[cat, below=of root] (jurint) {Jurnal\\Internasional};
\node[cat, below right=of root, xshift=10mm] (kon) {Konferensi};

\node[ven, below=of jur] (j1) {JSI, JUTI};
\node[ven, below=of j1] (j2) {JTIIK,\\TEKNOSI};
\node[ven, below=of j2] (j3) {JATISI,\\RESTI};

\node[ven, below=of jurint] (i1) {IJID};
\node[ven, below=of i1] (i2) {Jurnal Q3/Q4\\sistem informasi};

\node[ven, below=of kon] (k1) {ICITSI,\\ICACSIS};
\node[ven, below=of k1] (k2) {ICOIACT,\\ICITISEE};
\node[ven, below=of k2] (k3) {ICEEIE,\\ICoICT};

\draw (root) -- (jur);
\draw (root) -- (jurint);
\draw (root) -- (kon);
\draw (jur) -- (j1);
\draw (j1) -- (j2);
\draw (j2) -- (j3);
\draw (jurint) -- (i1);
\draw (i1) -- (i2);
\draw (kon) -- (k1);
\draw (k1) -- (k2);
\draw (k2) -- (k3);
\end{tikzpicture}
}
\caption{Tiga jenis target publikasi yang umum diakses oleh mahasiswa pascasarjana
Indonesia. Daftar tidak eksklusif; status akreditasi/indeks wajib diverifikasi
saat semester berjalan.}
\label{fig:venue-map}
\end{figure}

\textbf{Jurnal nasional terakreditasi (SINTA 2--4)} cocok untuk makalah dengan
kontribusi yang jelas dan analisis mendalam pada studi kasus Indonesia. Proses
telaah biasanya berlangsung beberapa bulan; tidak selalu ada tenggat kaku,
tetapi mahasiswa tetap perlu memperhatikan ritme penerbitan jurnal.

\textbf{Jurnal internasional (Q3--Q4)} cocok jika mahasiswa ingin pengalaman
proses telaah yang lebih ketat dan jangkauan pembaca yang lebih luas. Proses
telaah dapat berlangsung lebih lama, dan biaya publikasi pada sebagian jurnal
cukup besar. Reputasi jurnal harus selalu diperiksa agar mahasiswa terhindar
dari jurnal yang dicurigai \emph{predatory}.

\textbf{Konferensi nasional/internasional} cocok untuk mahasiswa yang ingin
menyelesaikan makalah dengan tenggat yang lebih jelas. Konferensi seperti
ICITSI, ICACSIS, ICOIACT, ICITISEE, ICEEIE, dan ICoICT dapat menjadi pintu
masuk pertama mahasiswa pascasarjana ke komunitas akademik yang lebih luas.
Status indeks, penerbit, jadwal, dan jalur/topik konferensi tetap harus
diverifikasi pada semester berjalan.

\subsection{Kriteria Pemilihan Target Publikasi}

Lima kriteria praktis berikut dapat membantu mahasiswa memilih target
publikasi:

\begin{enumerate}
  \item \textbf{Kesesuaian ruang lingkup:} apakah topik makalah cocok dengan
        ruang lingkup atau jalur/topik target publikasi?
  \item \textbf{Kelayakan akses data:} apakah data dan akses organisasi yang
        dimiliki cukup untuk memenuhi kedalaman analisis yang diharapkan?
  \item \textbf{Reputasi dan keberlanjutan:} apakah target publikasi memiliki
        reputasi yang jelas, terindeks dengan benar, dan terbit secara
        konsisten?
  \item \textbf{Tenggat dan kalender:} apakah tenggat pengajuan cocok dengan
        kalender akademik mata kuliah?
  \item \textbf{Biaya:} berapa biaya \emph{article processing charge} (APC)
        atau \emph{registration fee}? Apakah terjangkau atau tersedia subsidi?
\end{enumerate}

Pada Sesi~1, mahasiswa diharapkan memilih dua atau tiga target publikasi dan
melakukan verifikasi singkat terhadap kelima kriteria di atas. Verifikasi ini
bukan sekadar formalitas; ia akan memengaruhi struktur makalah yang ditulis
sepanjang semester.

\section{Tipe Makalah yang Dapat Dipilih}

Makalah yang dihasilkan dari mata kuliah ini dapat mengambil salah satu dari
empat bentuk berikut, sesuai minat mahasiswa dan akses data yang tersedia:

\begin{enumerate}
  \item \textbf{Studi kasus arsitektur enterprise:} analisis kondisi saat
        ini, rancangan target, dan peta jalan transformasi pada satu
        organisasi. Cocok jika mahasiswa
        memiliki akses cukup pada dokumen, proses, atau narasumber organisasi.
  \item \textbf{Design science atau rancangan solusi:} perancangan blueprint,
        strategi API, atau model tata kelola sebagai artefak solusi. Cocok jika
        masalah organisasi cukup jelas dan mahasiswa ingin menonjolkan
        kontribusi rancangan.
  \item \textbf{Evaluasi strategi platform/API:} analisis ekonomi API,
        ekosistem platform, network effects, dan tata kelola pada satu kasus.
        Cocok jika studi kasus memiliki aspek ekosistem, mitra, atau open API.
  \item \textbf{Kajian konseptual berbasis kasus lokal:} sintesis literatur dan
        contoh Indonesia untuk menghasilkan kerangka atau rekomendasi. Cocok
        jika akses data organisasi terbatas namun literatur dan data sekunder
        tersedia.
\end{enumerate}

Tipe makalah dipilih bersamaan dengan target publikasi di Sesi~1, lalu
diperhalus kembali pada Sesi~2--3 setelah rumusan masalah dan akses data lebih
jelas.

\section{Glosarium Mini}

Beberapa istilah kunci yang akan sering muncul:

\begin{itemize}
  \item \textbf{Enterprise Architecture (EA):} disiplin yang menyelaraskan
        strategi organisasi dengan kapabilitas teknologi melalui pemetaan,
        prinsip, dan keputusan terdokumentasi.
  \item \textbf{TOGAF-lite:} versi sederhana dari The Open Group Architecture
        Framework yang berfokus pada siklus inti: visi, kondisi saat ini,
        target, kesenjangan, dan peta jalan, tanpa membawa seluruh
        kompleksitas TOGAF.
  \item \textbf{BDAT:} singkatan dari Business, Data, Application, Technology
        sebagai empat domain arsitektur klasik.
  \item \textbf{ArchiMate:} bahasa pemodelan standar untuk arsitektur
        enterprise; menyusun arsitektur ke dalam lapisan motivasi, strategi,
        bisnis, aplikasi, teknologi/fisik, serta implementasi dan migrasi.
  \item \textbf{API (Application Programming Interface):} kontrak digital
        terstandar yang memungkinkan satu sistem berinteraksi dengan sistem
        lain.
  \item \textbf{Platform multi-sided:} model bisnis yang menghubungkan dua atau
        lebih kelompok pengguna dan menghasilkan \emph{network effect}.
  \item \textbf{Network effect:} fenomena bahwa nilai sebuah produk/platform
        meningkat bersamaan dengan bertambahnya pengguna pada satu atau lebih
        sisi platform.
  \item \textbf{Research question (RQ):} pertanyaan penelitian yang menjadi
        arah utama makalah dan dijawab melalui analisis serta pembahasan.
\end{itemize}

\section{Integrasi ke Makalah Publikasi}

Bab ini menyumbang bahan awal untuk dua bagian makalah:

\begin{itemize}
  \item \textbf{Pendahuluan:} pembaca akan menggunakan skenario organisasi
        sendiri (serupa dengan Bank Wanua) sebagai latar belakang makalah,
        kemudian
        menuliskan motivasi mengapa transformasi digital pada organisasi
        tersebut layak dianalisis.
  \item \textbf{Pernyataan kontribusi:} pernyataan singkat 2--3 kalimat
        tentang apa yang akan dihasilkan makalah dan mengapa hasil tersebut
        bernilai bagi komunitas akademik atau praktik.
\end{itemize}

Pada Sesi~2, RQ dan pernyataan kontribusi akan dipertajam. Pada Sesi~1, yang
dibutuhkan adalah draf awal yang cukup jelas untuk menjadi titik pijak.
Naskah belum perlu sempurna; yang penting cukup konkret untuk ditinjau bersama
dosen dan rekan sejawat.

\section{Diskusi Kritis dan Perangkap Umum}

Beberapa perangkap berikut sering muncul ketika mahasiswa pertama kali masuk
ke materi Enterprise Architecture:

\begin{enumerate}
  \item \emph{Memulai dari teknologi.} Mahasiswa sering tergoda memilih
        teknologi terlebih dahulu, misalnya ``saya ingin menulis tentang
        microservices'', alih-alih memulai dari masalah organisasi.
  \item \emph{Topik terlalu luas.} ``Transformasi digital BUMN energi
        nasional'' adalah judul tesis lima tahun, bukan makalah satu semester.
        Pilih satu organisasi, satu fokus.
  \item \emph{Mengabaikan konteks Indonesia.} Mengutip teori platform Silicon
        Valley tanpa menyelaraskan dengan UU PDP, OJK, atau realitas pasar
        lokal akan menghasilkan rekomendasi yang tidak dapat
        diimplementasikan.
  \item \emph{Memilih target publikasi setelah makalah selesai.} Akibatnya
        format, panjang, dan referensi tidak selaras dengan templat. Pilih
        target publikasi di Sesi~1 dan ikuti templat tersebut sejak awal.
  \item \emph{Menganggap blueprint sebagai luaran utama.} Blueprint adalah
        bukti analisis; luaran utama tetap makalah.
\end{enumerate}

\section{Latihan Terapan Individual}

Latihan pada bab ini sengaja dibuat ringan karena fungsinya adalah meletakkan
fondasi yang akan ditindaklanjuti pada bab-bab selanjutnya.

\subsection*{Latihan~2.1: Pemilihan Organisasi Studi Kasus (45 menit)}

Pilihlah satu organisasi nyata atau hipotetis-realistis dalam konteks
Indonesia yang akan Anda gunakan sebagai studi kasus sepanjang semester.
Tuliskan jawaban atas pertanyaan berikut secara ringkas (maksimum satu
halaman):

\begin{enumerate}
  \item Nama organisasi (boleh anonim/disamarkan, misalnya ``Bank Wanua'').
  \item Sektor (publik, pendidikan, kesehatan, keuangan, retail, logistik,
        manufaktur, sosial, layanan digital).
  \item Skala dan jangkauan operasi singkat.
  \item Mengapa organisasi ini menarik untuk dianalisis (1--2 paragraf).
  \item Akses data yang Anda miliki: dokumen publik, wawancara, observasi,
        data internal dengan izin, atau kombinasi.
\end{enumerate}

\subsection*{Latihan~2.2: Pernyataan Topik Awal (30 menit)}

Tuliskan satu paragraf (4--6 kalimat) sebagai pernyataan topik awal makalah
Anda. Paragraf ini sebaiknya menjawab tiga hal:

\begin{itemize}
  \item Apa masalah/peluang transformasi digital yang akan dianalisis?
  \item Pendekatan analisis apa yang akan digunakan (TOGAF-lite, design
        science, evaluasi platform, atau kajian konseptual)?
  \item Apa potensi kontribusi makalah bagi komunitas akademik atau praktik?
\end{itemize}

\subsection*{Latihan~2.3: Identifikasi 2--3 Target Publikasi (60 menit)}

Pilih dua atau tiga target publikasi (jurnal SINTA, jurnal internasional, atau
konferensi) yang relevan dengan topik Anda. Buat tabel ringkas dengan kolom:
nama target publikasi, ruang lingkup atau jalur/topik yang relevan, status
indeks atau akreditasi, biaya publikasi/APC, tenggat atau jadwal terbit, dan
tautan templat.

\subsection*{Latihan~2.4: Telaah Singkat Satu Makalah Contoh (60 menit)}

Cari satu makalah akademik (jurnal SINTA, jurnal internasional, atau prosiding
konferensi) yang relevan dengan topik Anda. Tuliskan refleksi satu halaman:

\begin{itemize}
  \item Apa pertanyaan penelitian (research question) makalah tersebut?
  \item Pendekatan/metode apa yang digunakan?
  \item Apa kontribusinya?
  \item Bagaimana makalah itu menempatkan konteks (Indonesia atau lainnya)?
  \item Apa yang dapat Anda pelajari dari makalah itu untuk naskah Anda?
\end{itemize}

\subsection*{Pertanyaan Refleksi}

\begin{enumerate}
  \item Apakah organisasi yang Anda pilih cukup spesifik untuk dianalisis
        dalam satu semester?
  \item Apakah akses data yang Anda miliki memungkinkan tipe makalah yang
        Anda pilih?
  \item Target publikasi mana yang paling realistis dilihat dari tenggat,
        biaya, dan kedalaman analisis yang dapat Anda capai?
  \item Setelah menelaah satu makalah contoh, bagian mana yang menurut Anda
        akan paling menantang ditulis?
\end{enumerate}

\section{Uji Mandiri}

\subsection*{Pertanyaan Konseptual}

\begin{enumerate}
  \item Jelaskan dengan kata-kata Anda sendiri mengapa Enterprise Architecture
        penting bagi pengambil keputusan non-IT.
  \item Sebutkan empat domain BDAT dan beri satu contoh artefak per domain.
  \item Apa yang dimaksud dengan \emph{network effect} pada platform
        multi-sided? Beri satu contoh dari Indonesia.
  \item Sebutkan minimal lima bagian pada anatomi makalah ilmiah dan jelaskan
        fungsi komunikasi setiap bagian.
  \item Apa perbedaan utama antara API \emph{private}, \emph{partner}, dan
        \emph{public}?
\end{enumerate}

\subsection*{Pertanyaan Aplikatif}

\begin{enumerate}
  \item Sebuah rumah sakit daerah ingin membangun layanan telekonsultasi.
        Sebutkan tiga keputusan EA strategis (lintas BDAT) yang perlu
        dipertimbangkan oleh direksi rumah sakit.
  \item Sebuah koperasi simpan-pinjam ingin membuka API agar UMKM dapat
        mengintegrasikan layanan pembayarannya. Apa risiko utama dari sisi
        regulasi dan tata kelola yang perlu diantisipasi?
\end{enumerate}

\subsection*{Pertanyaan Reflektif untuk Makalah}

\begin{enumerate}
  \item Bagian mana dari latihan~2.1--2.4 yang paling akan memengaruhi
        struktur Pendahuluan makalah Anda?
  \item Apakah ada kemungkinan organisasi studi kasus Anda berubah dalam dua
        bab pertama? Apa indikatornya?
\end{enumerate}

\section{Ringkasan Bab}

Bab ini memperkenalkan empat pilar yang menjadi dasar seluruh mata kuliah.
Pertama, Enterprise Architecture diposisikan sebagai jembatan dua arah antara
strategi organisasi dan kapabilitas teknologi. Kedua, kerangka BDAT (Business,
Data, Application, Technology) diperkenalkan sebagai cara memecah kerumitan
organisasi menjadi empat domain yang saling terkait. Ketiga, API dan platform
diposisikan sebagai pengungkit nilai pada model bisnis digital modern, dengan
studi kasus Indonesia sebagai sumber pembelajaran utama. Keempat, anatomi
makalah ilmiah diperkenalkan sebagai kerangka komunikasi yang akan dibangun
sesi demi sesi.

Konteks Indonesia dipertegas sebagai bagian inti analisis melalui regulasi
seperti UU PDP, UU ITE, PP 71/2019, Permenkominfo PSE, OJK, BSSN, dan SNAP
BI, serta melalui ekosistem digital yang akan sering digunakan sebagai studi
kasus. Pemilihan target publikasi juga diperkenalkan sebagai keputusan
strategis yang memengaruhi struktur makalah. Lima kriteria praktisnya adalah
kesesuaian ruang lingkup, kelayakan akses data, reputasi target publikasi,
tenggat, dan biaya.

Pada akhir bab, mahasiswa telah memilih satu organisasi studi kasus, menulis
pernyataan topik awal, dan mengidentifikasi dua sampai tiga kandidat target
publikasi. Tiga keluaran awal ini menjadi titik pijak Bab~3, ketika rumusan
masalah, pertanyaan penelitian, dan pernyataan kontribusi dipertajam melalui
\emph{lapisan motivasi} ArchiMate---pemangku kepentingan, pendorong, sasaran,
kebutuhan, dan kendala---pada studi kasus API Transfer Dana.

\section{Jembatan ke Bab Berikutnya}

Bab~3 masuk ke \emph{lapisan motivasi}---lapisan teratas dalam bahasa pemodelan
ArchiMate yang merekam \emph{mengapa} sebuah arsitektur dibangun. Untuk studi
kasus berjalan, yaitu API Transfer dan Perpindahan Dana Bank Wanua
(Bagian~\ref{sec:studi-kasus-berjalan}), pembaca memetakan pemangku kepentingan
beserta kekhawatirannya, pendorong perubahan, asesmen kondisi saat ini, sasaran,
hasil, kebutuhan, kendala, dan prinsip. Dari peta motivasi itu, ketidakselarasan
antara sasaran dan kapabilitas yang tersedia diubah menjadi rumusan masalah,
pertanyaan penelitian, dan pernyataan kontribusi makalah secara lebih tajam.

\section{Bacaan Lanjutan}

\subsection*{Bacaan Inti}

\begin{itemize}
  \item \emph{TOGAF Standard, 10th Edition}~\cite{togaf10}. Bagian pengantar
        (\emph{Architecture Development Method overview}) memberi gambaran
        besar sebelum pembaca masuk ke versi ``lite'' pada Bab~4.
  \item \emph{Designed for Digital}~\cite{ross2019designed}. Cocok untuk
        membangun intuisi tentang mengapa transformasi digital membutuhkan kerangka
        arsitektur.
  \item \emph{Platform Revolution}~\cite{parker2016platform}. Bab pertama
        mengenai pergeseran dari model linear ke platform sangat membantu
        untuk memahami kasus seperti Gojek dan Tokopedia.
\end{itemize}

\subsection*{Bacaan untuk Memperdalam}

\begin{itemize}
  \item \emph{APIs: A Strategy Guide}~\cite{jacobson2011apis}. Untuk pembaca
        yang ingin memperluas wawasan API sebelum Bab~10.
  \item \emph{The Business of Platforms}~\cite{cusumano2019business}.
        Bacaan pelengkap untuk memahami strategi dan tata kelola platform.
  \item \emph{Enterprise Architecture as Strategy}~\cite{ross2006enterprise}.
        Rujukan klasik EA yang relevan untuk argumen keselarasan bisnis--TI.
\end{itemize}

\subsection*{Sumber Indonesia}

\begin{itemize}
  \item \emph{Standar Nasional Open API Pembayaran (SNAP BI)}~\cite{snapbi}.
        Dokumen resmi bagi mahasiswa yang memilih topik perbankan/fintech.
  \item \emph{Satu Data Indonesia}~\cite{satudata}. Kerangka tata kelola data
        publik nasional yang relevan untuk studi kasus pemerintahan.
  \item \emph{Undang-Undang Pelindungan Data Pribadi}~\cite{uupdp2022} dan
        \emph{PP No.~71/2019}~\cite{pp712019} menjadi dua rujukan regulasi
        yang akan sering muncul sepanjang buku.
  \item Jurnal Sistem Informasi (JSI), JUTI, JTIIK, TEKNOSI, dan jurnal SINTA
        lain dengan ruang lingkup EA, sistem informasi, atau platform; gunakan SINTA
        \footnote{\url{https://sinta.kemdiktisaintek.go.id/}} dan Garuda
        \footnote{\url{https://garuda.kemdiktisaintek.go.id/}} untuk mencari
        makalah relevan.
\end{itemize}
